\section{Rangkuman}
\label{sec:rangkuman-bab14}

Pada bab ini telah dibahas:

\begin{itemize}
    \item Stack adalah daerah memori dengan aturan \eng{last in,
          first out}. Stack 8086 ditunjuk \reg{SS}:\reg{SP} dan
          tumbuh ke alamat yang lebih rendah.

    \item \asm{PUSH} mengurangi \reg{SP} dua dan menyimpan satu word.
          \asm{POP} mengambil satu word dan menambah \reg{SP} dua.

    \item Prosedur didefinisikan dengan \asm{PROC} dan \asm{ENDP}.
          Pemanggil memakai \asm{CALL}, sedangkan prosedur kembali
          dengan \asm{RET}.

    \item \asm{CALL} menyimpan alamat kembali ke stack. \asm{RET}
          mengambil alamat itu dan memuatnya kembali ke \reg{IP}.

    \item \eng{Stack frame} dibentuk dengan \asm{PUSH BP} dan
          \asm{MOV BP, SP}. Parameter pertama berada pada
          \code{[BP+4]}, parameter kedua pada \code{[BP+6]}.

    \item Data dapat disampaikan ke prosedur melalui register,
          \eng{stack}, atau variabel. \asm{RET} \code{n} membuang
          $n$ byte parameter setelah kembali.

    \item Macro didefinisikan dengan \asm{MACRO} dan diakhiri
          \asm{ENDM}. Macro dapat memiliki parameter dan disalin ke
          tempat pemakaian pada saat perakitan.

    \item Prosedur menekankan ukuran program dan modularitas,
          sedangkan macro menekankan kemudahan parameter dan
          menghilangkan biaya \asm{CALL}/\asm{RET}.
\end{itemize}

\begin{checklistbox}
\vspace{1.1em}
Tandai setiap butir yang sudah Anda kuasai.

\begin{itemize}
    \item[\checkbox] Saya dapat menjelaskan sifat LIFO pada stack.
    \item[\checkbox] Saya dapat menggunakan \asm{PUSH} dan
          \asm{POP}.
    \item[\checkbox] Saya dapat menjelaskan perubahan \reg{SP} pada
          operasi stack.
    \item[\checkbox] Saya dapat menulis prosedur dengan
          \asm{PROC}/\asm{ENDP}.
    \item[\checkbox] Saya dapat menggunakan \asm{CALL} dan
          \asm{RET}.
    \item[\checkbox] Saya dapat mengakses parameter melalui
          \eng{stack frame} dan \reg{BP}.
    \item[\checkbox] Saya dapat mendefinisikan dan memakai macro
          berparameter.
    \item[\checkbox] Saya dapat membedakan prosedur dan macro.
\end{itemize}

\vspace{0.4em}
Bila masih ada butir yang belum bertanda, ulangi bagian yang
bersesuaian sebelum melanjutkan ke bab berikutnya.
\end{checklistbox}

\begin{infobox}
\vspace{1.1em}
\textbf{Rujukan bab.} Materi pada bab ini diringkas dari
\cite{gunawan2019}, Bagian 12--13, hlm.~115--126, dan
\cite{sto2001}, BAB XIV, hlm.~73--77, serta BAB XVI--XVII,
hlm.~85--92.
\end{infobox}
